\section{Counting Patterns}
\label{sec:cses-counting-patterns}

\problemstatement{Given a text $t$ and $k$ patterns, count for each
pattern how many times it occurs in $t$ (occurrences may overlap).}

\begin{patternbox}
This is Finding Patterns with counts: the number of occurrences of a
substring equals the size of its \textbf{endpos} set, which the suffix
automaton stores per state. Compute those sizes once, then each pattern's
count is a lookup after walking to its state.
\end{patternbox}

\subsection*{Endpos sizes count occurrences}

In a suffix automaton, each state represents a set of substrings that
share the same set of ending positions (their \emph{endpos} set). The
number of occurrences of any substring in a state equals that state's
endpos-set size. To compute these: mark each state created as a genuine
prefix-endpoint (the ``\textit{last}'' after each extension) with count
$1$ and clones with $0$, then propagate counts up the suffix-link tree in
order of decreasing length -- each state adds its count to its link. A
pattern's count is the endpos size of the state its walk ends in.

\begin{ideabox}[Propagate Counts Up the Link Tree]
The suffix-link tree groups substrings by endpos containment: a child's
occurrences are a subset that also count toward the parent. Summing counts
from longer states into their links (processing by decreasing $\textit{len}$)
computes every endpos size in one linear pass.
\end{ideabox}

\subsection*{Algorithm}

\begin{enumerate}
  \item Build the SAM of $t$; set $\textit{cnt}=1$ for each non-clone
        state (the \textit{last} after each extend), $0$ for clones.
  \item Sort states by $\textit{len}$ descending; add each state's
        $\textit{cnt}$ to its suffix link's $\textit{cnt}$.
  \item For each pattern, walk to its state; output that state's
        $\textit{cnt}$ (or $0$ if the walk fails).
\end{enumerate}

\subsection*{Dry run}

\dryrun{$t=$``aaa''; pattern ``a''.}

\begin{itemize}
  \item ``a'' occurs at positions $1,2,3$ -- endpos size $3$.
  \item The SAM state for ``a'' accumulates count $3$ after propagation,
        which is the reported answer.
\end{itemize}

\subsection*{C++ implementation}

\begin{lstlisting}
#include <bits/stdc++.h>
using namespace std;

struct SAM {
    struct State { int len, link; map<char,int> next; };
    vector<State> st; vector<long long> cnt; int last;
    SAM() { st.push_back({0, -1, {}}); cnt.push_back(0); last = 0; }
    void extend(char c) {
        int cur = st.size();
        st.push_back({st[last].len + 1, -1, {}}); cnt.push_back(1);  // real endpoint
        int p = last;
        while (p != -1 && !st[p].next.count(c)) { st[p].next[c] = cur; p = st[p].link; }
        if (p == -1) st[cur].link = 0;
        else {
            int q = st[p].next[c];
            if (st[p].len + 1 == st[q].len) st[cur].link = q;
            else {
                int clone = st.size();
                st.push_back(st[q]); cnt.push_back(0);              // clone: no endpoint
                st[clone].len = st[p].len + 1;
                while (p != -1 && st[p].next[c] == q) { st[p].next[c] = clone; p = st[p].link; }
                st[q].link = st[cur].link = clone;
            }
        }
        last = cur;
    }
};

int main() {
    string t; int k;
    cin >> t >> k;
    SAM sam;
    for (char c : t) sam.extend(c);

    int m = sam.st.size();
    vector<int> order(m);
    for (int i = 0; i < m; i++) order[i] = i;
    sort(order.begin(), order.end(),
         [&](int a, int b){ return sam.st[a].len > sam.st[b].len; });
    for (int v : order)
        if (sam.st[v].link >= 0) sam.cnt[sam.st[v].link] += sam.cnt[v];

    while (k--) {
        string p; cin >> p;
        int cur = 0; bool ok = true;
        for (char c : p) {
            auto it = sam.st[cur].next.find(c);
            if (it == sam.st[cur].next.end()) { ok = false; break; }
            cur = it->second;
        }
        cout << (ok ? sam.cnt[cur] : 0) << "\n";
    }
    return 0;
}
\end{lstlisting}

\begin{complexitybox}
\textbf{Time Complexity.} $O(|t|\log\Sigma)$ build and count propagation,
plus $O(\sum|p_i|\log\Sigma)$ queries. \textbf{Space Complexity.}
$O(|t|)$.
\end{complexitybox}

\begin{takeawaybox}
Occurrence count = endpos-set size, computed by propagating counts up the
suffix-link tree. This single augmentation upgrades the SAM from ``does it
occur?'' to ``how many times?'' -- and the same endpos machinery drives
distinct-substring counting and ranking later in the chapter.
\end{takeawaybox}
